%% Literature Review - owner: Taj (taken over from Satya). 1 to 1.5 pages.
%% Every citation below refers to a source whose claims were checked against
%% the source itself, not paraphrased from memory. See docs/refs.bib for what
%% was verified and how.

This section traces the line of work our design sits on: the classification that
names it, the machines that first exploited it, its arrival in commodity
processors, and its present form on reconfigurable hardware.

\subsection{Flynn's Taxonomy}
Flynn~\cite{flynn1966} classified computing systems by the multiplicity of their
instruction and data streams, giving the four categories --- SISD, SIMD, MISD
and MIMD --- that the field still uses. The scheme separates two things easily conflated: issuing many
instructions at once, versus issuing one instruction that acts on many data
items. Our unit is squarely \emph{single instruction, multiple data} --- one
\texttt{VADD} is fetched and decoded once while four lanes act on four element
pairs. The scalar baseline in Section~\ref{sec:evaluation} is its SISD
counterpart, so comparing them measures what moving one category buys.

\subsection{Vector Processors}
The CRAY-1~\cite{russell1978} established the vector register machine. It
provided eight vector registers of sixty-four 64-bit elements each, a vector
length register, and a vector mask register for operating on individual
elements. Critically, it obtained its throughput by \emph{pipelining} functional
units deeply rather than by replicating them, and by \emph{chaining} --- feeding
the result of one vector instruction directly into the next without a round trip
through memory.

We made the opposite choice: our unit replicates the ALU four times and
completes an operation in one cycle, where the CRAY-1 streamed elements through
one deep pipeline. Replication costs area linearly and keeps the critical path
short; pipelining costs latency and control complexity but scales to far longer
vectors. At four lanes and a 16-element benchmark, replication is the cheaper
choice, and Section~\ref{sec:evaluation} quantifies it in LUTs. Hennessy and
Patterson~\cite{hennessy2017} treat the two as the axes along which vector
hardware scales.

\subsection{Short-Vector SIMD in Commodity Processors}
The idea reached general-purpose CPUs in a narrower form. Extensions such as
SSE, AVX and ARM NEON add fixed-width registers treated as packed collections of
smaller elements. The width is fixed by the encoding, so code compiled for one
width cannot exploit a wider machine, and a new extension is needed each time
the register file grows. Our format shares this property: \texttt{VADD} names
128-bit registers, and widening the datapath changes what the encoding means.

The RISC-V vector extension~\cite{riscvvector} was designed to avoid exactly
that. Its programming model is vector-length agnostic: software queries the
hardware for the vector length it can support and loops accordingly, so the same
binary runs on implementations with different amounts of hardware. This is the
clearest contrast with what we built. Our \texttt{LANES} parameter is a
\emph{synthesis-time} decision that changes the instruction set's meaning,
whereas the RISC-V approach makes vector length a \emph{run-time} property the
program adapts to. The parameterisation in Section~\ref{sec:datapath} buys
configurability for the hardware designer, not for the compiled program.

\subsection{Soft Vector Processors on FPGAs}
The work closest to ours is Yu, Lemieux and Eagleston~\cite{yu2008soft}, who
proposed adding a vector core to an FPGA soft processor as a general-purpose
accelerator, instead of designing a custom accelerator for each application.
Their argument is development cost: a programmer with no hardware experience can
trade area for performance by scaling the lane count. Across three benchmark
kernels they report speedups of roughly 3--29$\times$ for 6--30$\times$ the
area of a Nios~II/s processor.

Their lane-scaling knob is the direct ancestor of our \texttt{LANES} parameter,
and their area-versus-performance framing is the one we adopt in
Section~\ref{sec:evaluation}, where we report a measured $4.00\times$ speed-up
against a $4.03\times$ area cost. The ratio is far more favourable than theirs,
but the comparison is not like-for-like and should not be read as an
improvement: our benchmark is a single, perfectly data-parallel kernel on an
element count that is an exact multiple of the lane width, whereas their figures
span real kernels with imperfect vectorisation. The honest conclusion is that
our result confirms the expected linear trade under ideal conditions, which is
the regime in which their lower bound sits.
